%%%%%%%%%%%%%%%%%%%%%%%%%%%%%%%%%%%%%%%%%%%%%%%%%%%%%%%%%%%%
% 7.3 DISCRETE DISTRIBUTIONS
% The Composition of Advanced Mathematics
%%%%%%%%%%%%%%%%%%%%%%%%%%%%%%%%%%%%%%%%%%%%%%%%%%%%%%%%%%%%

\section{Discrete Distributions}
\label{sec:discrete-distributions}

\prerequisite{
Elementary probability, random variables, probability mass functions,
cumulative distribution functions, combinations, conditional
probability, independence, and basic infinite series.
}

\learningobjectives{
By the end of this section, the reader should be able to:
\begin{itemize}
    \item recognize and construct discrete probability distributions;
    \item verify that a proposed probability mass function is valid;
    \item identify Bernoulli random variables;
    \item use the binomial distribution for repeated independent trials;
    \item use the geometric distribution for waiting-time problems;
    \item use the negative binomial distribution;
    \item distinguish geometric and negative binomial conventions;
    \item use the hypergeometric distribution for sampling without
          replacement;
    \item distinguish binomial and hypergeometric models;
    \item use the Poisson distribution for event counts;
    \item understand the Poisson approximation to the binomial;
    \item use the discrete uniform distribution;
    \item work with multinomial distributions;
    \item recognize occupancy and categorical counting models;
    \item use probability-generating functions at an introductory level;
    \item understand recurrence relations between neighboring
          probabilities;
    \item derive cumulative probabilities from PMFs;
    \item recognize limiting relationships among common discrete
          distributions;
    \item select an appropriate discrete model from a verbal problem;
    \item understand how standard distributions encode repeated
          probabilistic structures.
\end{itemize}
}

%%%%%%%%%%%%%%%%%%%%%%%%%%%%%%%%%%%%%%%%%%%%%%%%%%%%%%%%%%%%
\subsection{What Is a Discrete Distribution?}
\label{subsec:what-is-discrete-distribution}
%%%%%%%%%%%%%%%%%%%%%%%%%%%%%%%%%%%%%%%%%%%%%%%%%%%%%%%%%%%%

A discrete random variable takes values in a finite or countably
infinite set.

Its distribution is described by a probability mass function

\[
\boxed{
p_X(x)=P(X=x).
}
\]

A valid PMF must satisfy

\[
\boxed{
p_X(x)\geq0
}
\]

for every \(x\), and

\[
\boxed{
\sum_x p_X(x)=1.
}
\]

Standard discrete distributions arise when the same probabilistic
structures appear repeatedly.

%%%%%%%%%%%%%%%%%%%%%%%%%%%%%%%%%%%%%%%%%%%%%%%%%%%%%%%%%%%%
\subsection{Why Standard Distributions Matter}
\label{subsec:why-standard-discrete-distributions}
%%%%%%%%%%%%%%%%%%%%%%%%%%%%%%%%%%%%%%%%%%%%%%%%%%%%%%%%%%%%

Many probability problems share the same underlying structure.

For example:

\[
\boxed{
\text{one success/failure trial}
\longrightarrow
\text{Bernoulli},
}
\]

\[
\boxed{
\text{number of successes in \(n\) independent trials}
\longrightarrow
\text{Binomial},
}
\]

\[
\boxed{
\text{trials until first success}
\longrightarrow
\text{Geometric},
}
\]

\[
\boxed{
\text{sampling without replacement}
\longrightarrow
\text{Hypergeometric},
}
\]

and

\[
\boxed{
\text{counting rare events in a fixed interval}
\longrightarrow
\text{Poisson}.
}
\]

Recognizing the structure often determines the correct model immediately.

%%%%%%%%%%%%%%%%%%%%%%%%%%%%%%%%%%%%%%%%%%%%%%%%%%%%%%%%%%%%
\subsection{Discrete Uniform Distribution}
\label{subsec:discrete-uniform-distribution}
%%%%%%%%%%%%%%%%%%%%%%%%%%%%%%%%%%%%%%%%%%%%%%%%%%%%%%%%%%%%

\begin{definitionbox}{Discrete Uniform Distribution}
A random variable \(X\) is discrete uniform on a finite set

\[
\{x_1,\ldots,x_n\}
\]

if each value is equally likely:

\[
\boxed{
P(X=x_i)=\frac1n.
}
\]
\end{definitionbox}

One notation is

\[
X\sim\operatorname{Unif}\{x_1,\ldots,x_n\}.
\]

%%%%%%%%%%%%%%%%%%%%%%%%%%%%%%%%%%%%%%%%%%%%%%%%%%%%%%%%%%%%
\subsection{Worked Example: Fair Die}
\label{subsec:worked-discrete-uniform-die}
%%%%%%%%%%%%%%%%%%%%%%%%%%%%%%%%%%%%%%%%%%%%%%%%%%%%%%%%%%%%

\begin{workedexamplebox}{A Fair Six-Sided Die}

Let \(X\) be the number shown on a fair die.

Then

\[
X\in\{1,2,3,4,5,6\}
\]

and

\[
P(X=x)=\frac16
\]

for each

\[
x\in\{1,\ldots,6\}.
\]

\begin{answerbox}
\[
\boxed{
p_X(x)
=
\begin{cases}
\dfrac16,&x=1,2,3,4,5,6,\\[2mm]
0,&\text{otherwise}.
\end{cases}
}
\]
\end{answerbox}

\end{workedexamplebox}

%%%%%%%%%%%%%%%%%%%%%%%%%%%%%%%%%%%%%%%%%%%%%%%%%%%%%%%%%%%%
\subsection{Bernoulli Distribution}
\label{subsec:bernoulli-distribution}
%%%%%%%%%%%%%%%%%%%%%%%%%%%%%%%%%%%%%%%%%%%%%%%%%%%%%%%%%%%%

A Bernoulli random variable models one trial with two outcomes.

\begin{definitionbox}{Bernoulli Distribution}
A random variable \(X\) has a Bernoulli distribution with parameter
\(p\), written

\[
\boxed{
X\sim\operatorname{Bernoulli}(p),
}
\]

if

\[
\boxed{
P(X=1)=p
}
\]

and

\[
\boxed{
P(X=0)=1-p.
}
\]
\end{definitionbox}

Usually,

\[
1
\]

represents success and

\[
0
\]

represents failure.

%%%%%%%%%%%%%%%%%%%%%%%%%%%%%%%%%%%%%%%%%%%%%%%%%%%%%%%%%%%%
\subsection{Bernoulli PMF}
\label{subsec:bernoulli-pmf}
%%%%%%%%%%%%%%%%%%%%%%%%%%%%%%%%%%%%%%%%%%%%%%%%%%%%%%%%%%%%

The Bernoulli PMF can be written compactly as

\[
\boxed{
p_X(x)
=
p^x(1-p)^{1-x},
\qquad
x\in\{0,1\}.
}
\]

Indeed,

\[
x=1
\]

gives

\[
p_X(1)=p,
\]

while

\[
x=0
\]

gives

\[
p_X(0)=1-p.
\]

%%%%%%%%%%%%%%%%%%%%%%%%%%%%%%%%%%%%%%%%%%%%%%%%%%%%%%%%%%%%
\subsection{Indicator Variables Are Bernoulli}
\label{subsec:indicator-bernoulli}
%%%%%%%%%%%%%%%%%%%%%%%%%%%%%%%%%%%%%%%%%%%%%%%%%%%%%%%%%%%%

For an event \(A\), the indicator

\[
\mathbf{1}_A
\]

satisfies

\[
P(\mathbf{1}_A=1)=P(A).
\]

Therefore,

\[
\boxed{
\mathbf{1}_A
\sim
\operatorname{Bernoulli}(P(A)).
}
\]

This connects event probabilities directly with standard discrete
distributions.

%%%%%%%%%%%%%%%%%%%%%%%%%%%%%%%%%%%%%%%%%%%%%%%%%%%%%%%%%%%%
\subsection{Bernoulli Trials}
\label{subsec:bernoulli-trials}
%%%%%%%%%%%%%%%%%%%%%%%%%%%%%%%%%%%%%%%%%%%%%%%%%%%%%%%%%%%%

A sequence of Bernoulli trials consists of repeated experiments having:

\begin{itemize}
    \item exactly two outcomes;
    \item constant success probability \(p\);
    \item independent trials.
\end{itemize}

Repeated Bernoulli trials lead to several important distributions.

%%%%%%%%%%%%%%%%%%%%%%%%%%%%%%%%%%%%%%%%%%%%%%%%%%%%%%%%%%%%
\subsection{Binomial Distribution}
\label{subsec:binomial-distribution}
%%%%%%%%%%%%%%%%%%%%%%%%%%%%%%%%%%%%%%%%%%%%%%%%%%%%%%%%%%%%

\begin{definitionbox}{Binomial Distribution}
Suppose \(n\) independent Bernoulli trials are performed, each with
success probability \(p\).

Let

\[
X
=
\text{number of successes}.
\]

Then

\[
\boxed{
X\sim\operatorname{Bin}(n,p).
}
\]
\end{definitionbox}

The parameters are:

\[
n
=
\text{number of trials},
\]

and

\[
p
=
\text{success probability per trial}.
\]

%%%%%%%%%%%%%%%%%%%%%%%%%%%%%%%%%%%%%%%%%%%%%%%%%%%%%%%%%%%%
\subsection{Deriving the Binomial PMF}
\label{subsec:derive-binomial-pmf}
%%%%%%%%%%%%%%%%%%%%%%%%%%%%%%%%%%%%%%%%%%%%%%%%%%%%%%%%%%%%

A particular sequence containing exactly \(k\) successes and
\(n-k\) failures has probability

\[
p^k(1-p)^{n-k}.
\]

There are

\[
\binom nk
\]

ways to choose which trials are successes.

Therefore,

\[
\boxed{
P(X=k)
=
\binom nk
p^k(1-p)^{n-k},
}
\]

for

\[
k=0,1,\ldots,n.
\]

%%%%%%%%%%%%%%%%%%%%%%%%%%%%%%%%%%%%%%%%%%%%%%%%%%%%%%%%%%%%
\subsection{Binomial PMF}
\label{subsec:binomial-pmf}
%%%%%%%%%%%%%%%%%%%%%%%%%%%%%%%%%%%%%%%%%%%%%%%%%%%%%%%%%%%%

The full PMF is

\[
\boxed{
p_X(k)
=
\begin{cases}
\displaystyle
\binom nk
p^k(1-p)^{n-k},
&
k=0,1,\ldots,n,
\\[3mm]
0,
&
\text{otherwise}.
\end{cases}
}
\]

The support is

\[
\boxed{
\{0,1,\ldots,n\}.
}
\]

%%%%%%%%%%%%%%%%%%%%%%%%%%%%%%%%%%%%%%%%%%%%%%%%%%%%%%%%%%%%
\subsection{Why the Binomial PMF Sums to One}
\label{subsec:binomial-normalization}
%%%%%%%%%%%%%%%%%%%%%%%%%%%%%%%%%%%%%%%%%%%%%%%%%%%%%%%%%%%%

Using the binomial theorem,

\[
\sum_{k=0}^{n}
\binom nk
p^k(1-p)^{n-k}
=
(p+(1-p))^n.
\]

Therefore,

\[
\boxed{
\sum_{k=0}^{n}
P(X=k)
=
1^n
=
1.
}
\]

Thus the binomial PMF is properly normalized.

%%%%%%%%%%%%%%%%%%%%%%%%%%%%%%%%%%%%%%%%%%%%%%%%%%%%%%%%%%%%
\subsection{Worked Example: Exactly Four Successes}
\label{subsec:worked-binomial-four-successes}
%%%%%%%%%%%%%%%%%%%%%%%%%%%%%%%%%%%%%%%%%%%%%%%%%%%%%%%%%%%%

\begin{workedexamplebox}{Binomial Probability}

Suppose

\[
X\sim\operatorname{Bin}(10,0.3).
\]

Find

\[
P(X=4).
\]

Use

\[
P(X=4)
=
\binom{10}{4}
(0.3)^4
(0.7)^6.
\]

\begin{answerbox}
\[
\boxed{
P(X=4)
=
\binom{10}{4}
(0.3)^4
(0.7)^6.
}
\]
\end{answerbox}

\end{workedexamplebox}

%%%%%%%%%%%%%%%%%%%%%%%%%%%%%%%%%%%%%%%%%%%%%%%%%%%%%%%%%%%%
\subsection{At Most and At Least for Binomial Variables}
\label{subsec:binomial-cumulative-events}
%%%%%%%%%%%%%%%%%%%%%%%%%%%%%%%%%%%%%%%%%%%%%%%%%%%%%%%%%%%%

If

\[
X\sim\operatorname{Bin}(n,p),
\]

then

\[
P(X\leq k)
=
\sum_{j=0}^{k}
\binom njp^j(1-p)^{n-j}.
\]

Similarly,

\[
P(X\geq k)
=
\sum_{j=k}^{n}
\binom njp^j(1-p)^{n-j}.
\]

Often it is easier to use a complement:

\[
\boxed{
P(X\geq1)
=
1-P(X=0)
=
1-(1-p)^n.
}
\]

%%%%%%%%%%%%%%%%%%%%%%%%%%%%%%%%%%%%%%%%%%%%%%%%%%%%%%%%%%%%
\subsection{Worked Example: At Least One Success}
\label{subsec:worked-binomial-at-least-one}
%%%%%%%%%%%%%%%%%%%%%%%%%%%%%%%%%%%%%%%%%%%%%%%%%%%%%%%%%%%%

\begin{workedexamplebox}{At Least One Defective Item}

Suppose each component is defective independently with probability

\[
0.02.
\]

Inspect \(20\) components.

Let

\[
X\sim\operatorname{Bin}(20,0.02).
\]

Then

\[
P(X\geq1)
=
1-P(X=0).
\]

Therefore,

\[
\begin{answerbox}
\boxed{
P(X\geq1)
=
1-(0.98)^{20}.
}
\]
\end{answerbox}

\end{workedexamplebox}

%%%%%%%%%%%%%%%%%%%%%%%%%%%%%%%%%%%%%%%%%%%%%%%%%%%%%%%%%%%%
\subsection{When Is the Binomial Model Appropriate?}
\label{subsec:binomial-conditions}
%%%%%%%%%%%%%%%%%%%%%%%%%%%%%%%%%%%%%%%%%%%%%%%%%%%%%%%%%%%%

A binomial model is appropriate when:

\begin{enumerate}
    \item the number of trials \(n\) is fixed;
    \item each trial has two relevant outcomes;
    \item the trials are independent;
    \item each trial has the same success probability \(p\);
    \item \(X\) counts the number of successes.
\end{enumerate}

If one of these assumptions fails, another model may be more appropriate.

%%%%%%%%%%%%%%%%%%%%%%%%%%%%%%%%%%%%%%%%%%%%%%%%%%%%%%%%%%%%
\subsection{Binomial Recurrence Between Probabilities}
\label{subsec:binomial-probability-ratio}
%%%%%%%%%%%%%%%%%%%%%%%%%%%%%%%%%%%%%%%%%%%%%%%%%%%%%%%%%%%%

For

\[
X\sim\operatorname{Bin}(n,p),
\]

neighboring PMF values satisfy

\[
\frac{P(X=k+1)}{P(X=k)}
=
\frac{n-k}{k+1}
\frac{p}{1-p}.
\]

Therefore,

\[
\boxed{
P(X=k+1)
=
P(X=k)
\frac{n-k}{k+1}
\frac{p}{1-p}.
}
\]

This relation is useful for studying the shape and mode of the
distribution.

%%%%%%%%%%%%%%%%%%%%%%%%%%%%%%%%%%%%%%%%%%%%%%%%%%%%%%%%%%%%
\subsection{Mode of the Binomial Distribution}
\label{subsec:binomial-mode}
%%%%%%%%%%%%%%%%%%%%%%%%%%%%%%%%%%%%%%%%%%%%%%%%%%%%%%%%%%%%

The most likely value occurs near

\[
(n+1)p.
\]

More precisely, if

\[
(n+1)p
\]

is not an integer, the unique mode is

\[
\boxed{
\lfloor(n+1)p\rfloor.
}
\]

If \((n+1)p\) is an integer, there are two adjacent modes.

%%%%%%%%%%%%%%%%%%%%%%%%%%%%%%%%%%%%%%%%%%%%%%%%%%%%%%%%%%%%
\subsection{Geometric Distribution}
\label{subsec:geometric-distribution}
%%%%%%%%%%%%%%%%%%%%%%%%%%%%%%%%%%%%%%%%%%%%%%%%%%%%%%%%%%%%

The geometric distribution models the waiting time until the first
success in repeated independent Bernoulli trials.

There are two common conventions.

In this book, unless otherwise stated, \(X\) counts the number of trials
until the first success.

%%%%%%%%%%%%%%%%%%%%%%%%%%%%%%%%%%%%%%%%%%%%%%%%%%%%%%%%%%%%
\subsection{Geometric PMF}
\label{subsec:geometric-pmf}
%%%%%%%%%%%%%%%%%%%%%%%%%%%%%%%%%%%%%%%%%%%%%%%%%%%%%%%%%%%%

\begin{definitionbox}{Geometric Distribution}
Let independent Bernoulli trials have success probability \(p\).

If \(X\) is the number of trials until the first success, then

\[
\boxed{
X\sim\operatorname{Geom}(p)
}
\]

with PMF

\[
\boxed{
P(X=k)
=
(1-p)^{k-1}p,
\qquad
k=1,2,3,\ldots.
}
\]
\end{definitionbox}

The first \(k-1\) trials must fail and the \(k\)-th must succeed.

%%%%%%%%%%%%%%%%%%%%%%%%%%%%%%%%%%%%%%%%%%%%%%%%%%%%%%%%%%%%
\subsection{Geometric Normalization}
\label{subsec:geometric-normalization}
%%%%%%%%%%%%%%%%%%%%%%%%%%%%%%%%%%%%%%%%%%%%%%%%%%%%%%%%%%%%

We verify:

\[
\sum_{k=1}^{\infty}
(1-p)^{k-1}p
=
p
\sum_{j=0}^{\infty}(1-p)^j.
\]

Using the geometric-series formula,

\[
\sum_{j=0}^{\infty}(1-p)^j
=
\frac1p.
\]

Therefore,

\[
\boxed{
\sum_{k=1}^{\infty}
P(X=k)=1.
}
\]

%%%%%%%%%%%%%%%%%%%%%%%%%%%%%%%%%%%%%%%%%%%%%%%%%%%%%%%%%%%%
\subsection{Worked Example: First Success}
\label{subsec:worked-geometric}
%%%%%%%%%%%%%%%%%%%%%%%%%%%%%%%%%%%%%%%%%%%%%%%%%%%%%%%%%%%%

\begin{workedexamplebox}{First Success on Trial Five}

Suppose each trial succeeds independently with probability

\[
p=0.2.
\]

Find the probability that the first success occurs on trial \(5\).

The first four trials must fail and the fifth must succeed:

\[
P(X=5)
=
(0.8)^4(0.2).
\]

\begin{answerbox}
\[
\boxed{
P(X=5)
=
(0.8)^4(0.2).
}
\]
\end{answerbox}

\end{workedexamplebox}

%%%%%%%%%%%%%%%%%%%%%%%%%%%%%%%%%%%%%%%%%%%%%%%%%%%%%%%%%%%%
\subsection{Geometric Tail Probability}
\label{subsec:geometric-tail}
%%%%%%%%%%%%%%%%%%%%%%%%%%%%%%%%%%%%%%%%%%%%%%%%%%%%%%%%%%%%

For

\[
X\sim\operatorname{Geom}(p),
\]

the event

\[
X>k
\]

means the first \(k\) trials all fail.

Therefore,

\[
\boxed{
P(X>k)
=
(1-p)^k.
}
\]

Consequently,

\[
\boxed{
P(X\leq k)
=
1-(1-p)^k.
}
\]

%%%%%%%%%%%%%%%%%%%%%%%%%%%%%%%%%%%%%%%%%%%%%%%%%%%%%%%%%%%%
\subsection{Memoryless Property}
\label{subsec:geometric-memoryless}
%%%%%%%%%%%%%%%%%%%%%%%%%%%%%%%%%%%%%%%%%%%%%%%%%%%%%%%%%%%%

The geometric distribution is memoryless.

\begin{theorem}[Memoryless Property]
If

\[
X\sim\operatorname{Geom}(p),
\]

then

\[
\boxed{
P(X>s+t\mid X>s)
=
P(X>t).
}
\]
\end{theorem}

The fact that no success has occurred during the first \(s\) trials does
not change the distribution of the additional waiting time.

%%%%%%%%%%%%%%%%%%%%%%%%%%%%%%%%%%%%%%%%%%%%%%%%%%%%%%%%%%%%
\subsection{Proof of the Geometric Memoryless Property}
\label{subsec:proof-geometric-memoryless}
%%%%%%%%%%%%%%%%%%%%%%%%%%%%%%%%%%%%%%%%%%%%%%%%%%%%%%%%%%%%

Using the geometric tail,

\[
P(X>s+t\mid X>s)
=
\frac{P(X>s+t)}{P(X>s)}.
\]

Then

\[
=
\frac{(1-p)^{s+t}}{(1-p)^s}.
\]

Therefore,

\[
\boxed{
P(X>s+t\mid X>s)
=
(1-p)^t
=
P(X>t).
}
\]

%%%%%%%%%%%%%%%%%%%%%%%%%%%%%%%%%%%%%%%%%%%%%%%%%%%%%%%%%%%%
\subsection{Alternative Geometric Convention}
\label{subsec:alternative-geometric-convention}
%%%%%%%%%%%%%%%%%%%%%%%%%%%%%%%%%%%%%%%%%%%%%%%%%%%%%%%%%%%%

Some sources define a geometric random variable as the number of
failures before the first success.

Under that convention, the support is

\[
\{0,1,2,\ldots\}
\]

and

\[
\boxed{
P(Y=k)
=
(1-p)^kp.
}
\]

The two conventions satisfy

\[
\boxed{
X=Y+1.
}
\]

\begin{warningbox}
Always check whether a geometric variable counts trials until success
or failures before success.
\end{warningbox}

%%%%%%%%%%%%%%%%%%%%%%%%%%%%%%%%%%%%%%%%%%%%%%%%%%%%%%%%%%%%
\subsection{Negative Binomial Distribution}
\label{subsec:negative-binomial-distribution}
%%%%%%%%%%%%%%%%%%%%%%%%%%%%%%%%%%%%%%%%%%%%%%%%%%%%%%%%%%%%

The negative binomial distribution generalizes the geometric
distribution.

Instead of waiting for the first success, one waits for the \(r\)-th
success.

%%%%%%%%%%%%%%%%%%%%%%%%%%%%%%%%%%%%%%%%%%%%%%%%%%%%%%%%%%%%
\subsection{Negative Binomial PMF}
\label{subsec:negative-binomial-pmf}
%%%%%%%%%%%%%%%%%%%%%%%%%%%%%%%%%%%%%%%%%%%%%%%%%%%%%%%%%%%%

\begin{definitionbox}{Negative Binomial Distribution}
Suppose independent Bernoulli trials have success probability \(p\).

Let \(X\) be the number of trials required to obtain the \(r\)-th
success.

Then

\[
\boxed{
X\sim\operatorname{NegBin}(r,p)
}
\]

with

\[
\boxed{
P(X=k)
=
\binom{k-1}{r-1}
p^r
(1-p)^{k-r},
}
\]

for

\[
k=r,r+1,r+2,\ldots.
\]
\end{definitionbox}

%%%%%%%%%%%%%%%%%%%%%%%%%%%%%%%%%%%%%%%%%%%%%%%%%%%%%%%%%%%%
\subsection{Deriving the Negative Binomial PMF}
\label{subsec:derive-negative-binomial}
%%%%%%%%%%%%%%%%%%%%%%%%%%%%%%%%%%%%%%%%%%%%%%%%%%%%%%%%%%%%

For the \(r\)-th success to occur on trial \(k\):

\begin{itemize}
    \item trial \(k\) must be a success;
    \item among the first \(k-1\) trials, exactly \(r-1\) must be
          successes.
\end{itemize}

There are

\[
\binom{k-1}{r-1}
\]

ways to arrange those earlier successes.

Hence,

\[
\boxed{
P(X=k)
=
\binom{k-1}{r-1}
p^r(1-p)^{k-r}.
}
\]

%%%%%%%%%%%%%%%%%%%%%%%%%%%%%%%%%%%%%%%%%%%%%%%%%%%%%%%%%%%%
\subsection{Worked Example: Third Success}
\label{subsec:worked-negative-binomial}
%%%%%%%%%%%%%%%%%%%%%%%%%%%%%%%%%%%%%%%%%%%%%%%%%%%%%%%%%%%%

\begin{workedexamplebox}{Third Success on Trial Seven}

Suppose

\[
p=0.4.
\]

Find the probability that the third success occurs on trial \(7\).

Among the first \(6\) trials, exactly \(2\) must succeed.

The seventh trial must also succeed.

Thus,

\[
P(X=7)
=
\binom62
(0.4)^3
(0.6)^4.
\]

\begin{answerbox}
\[
\boxed{
P(X=7)
=
\binom62
(0.4)^3
(0.6)^4.
}
\]
\end{answerbox}

\end{workedexamplebox}

%%%%%%%%%%%%%%%%%%%%%%%%%%%%%%%%%%%%%%%%%%%%%%%%%%%%%%%%%%%%
\subsection{Geometric as a Special Case}
\label{subsec:geometric-negative-binomial}
%%%%%%%%%%%%%%%%%%%%%%%%%%%%%%%%%%%%%%%%%%%%%%%%%%%%%%%%%%%%

Set

\[
r=1.
\]

Then

\[
\binom{k-1}{0}=1,
\]

so the negative binomial PMF becomes

\[
(1-p)^{k-1}p.
\]

Therefore,

\[
\boxed{
\operatorname{Geom}(p)
=
\operatorname{NegBin}(1,p)
}
\]

under the trial-counting convention used here.

%%%%%%%%%%%%%%%%%%%%%%%%%%%%%%%%%%%%%%%%%%%%%%%%%%%%%%%%%%%%
\subsection{Hypergeometric Distribution}
\label{subsec:hypergeometric-distribution}
%%%%%%%%%%%%%%%%%%%%%%%%%%%%%%%%%%%%%%%%%%%%%%%%%%%%%%%%%%%%

The hypergeometric distribution models sampling without replacement
from a finite population.

Suppose:

\[
N
=
\text{population size},
\]

\[
K
=
\text{number of success-type objects},
\]

and

\[
n
=
\text{sample size}.
\]

Let \(X\) count the number of successes in the sample.

%%%%%%%%%%%%%%%%%%%%%%%%%%%%%%%%%%%%%%%%%%%%%%%%%%%%%%%%%%%%
\subsection{Hypergeometric PMF}
\label{subsec:hypergeometric-pmf}
%%%%%%%%%%%%%%%%%%%%%%%%%%%%%%%%%%%%%%%%%%%%%%%%%%%%%%%%%%%%

\begin{definitionbox}{Hypergeometric Distribution}
If \(n\) objects are sampled uniformly without replacement from a
population of size \(N\) containing \(K\) success-type objects, then

\[
\boxed{
X\sim\operatorname{Hypergeom}(N,K,n)
}
\]

with PMF

\[
\boxed{
P(X=k)
=
\frac{
\binom Kk
\binom{N-K}{n-k}
}{
\binom Nn
}.
}
\]
\end{definitionbox}

%%%%%%%%%%%%%%%%%%%%%%%%%%%%%%%%%%%%%%%%%%%%%%%%%%%%%%%%%%%%
\subsection{Hypergeometric Support}
\label{subsec:hypergeometric-support}
%%%%%%%%%%%%%%%%%%%%%%%%%%%%%%%%%%%%%%%%%%%%%%%%%%%%%%%%%%%%

The possible values must satisfy all counting constraints.

Therefore,

\[
\boxed{
\max(0,n-(N-K))
\leq
k
\leq
\min(n,K).
}
\]

This ensures there are enough successes and failures available to form
the sample.

%%%%%%%%%%%%%%%%%%%%%%%%%%%%%%%%%%%%%%%%%%%%%%%%%%%%%%%%%%%%
\subsection{Deriving the Hypergeometric PMF}
\label{subsec:derive-hypergeometric}
%%%%%%%%%%%%%%%%%%%%%%%%%%%%%%%%%%%%%%%%%%%%%%%%%%%%%%%%%%%%

There are

\[
\binom Nn
\]

possible samples of size \(n\).

To obtain exactly \(k\) successes:

\[
\binom Kk
\]

ways choose the successes, and

\[
\binom{N-K}{n-k}
\]

ways choose the failures.

Therefore,

\[
\boxed{
P(X=k)
=
\frac{
\binom Kk
\binom{N-K}{n-k}
}{
\binom Nn
}.
}
\]

%%%%%%%%%%%%%%%%%%%%%%%%%%%%%%%%%%%%%%%%%%%%%%%%%%%%%%%%%%%%
\subsection{Worked Example: Defective Components}
\label{subsec:worked-hypergeometric}
%%%%%%%%%%%%%%%%%%%%%%%%%%%%%%%%%%%%%%%%%%%%%%%%%%%%%%%%%%%%

\begin{workedexamplebox}{Sampling Without Replacement}

A box contains \(20\) components, of which \(5\) are defective.

Select \(4\) components without replacement.

Let \(X\) be the number of defective components selected.

Then

\[
X
\sim
\operatorname{Hypergeom}(20,5,4).
\]

The probability of exactly \(2\) defective components is

\[
P(X=2)
=
\frac{
\binom52
\binom{15}{2}
}{
\binom{20}{4}
}.
\]

\begin{answerbox}
\[
\boxed{
P(X=2)
=
\frac{
\binom52\binom{15}{2}
}{
\binom{20}{4}
}.
}
\]
\end{answerbox}

\end{workedexamplebox}

%%%%%%%%%%%%%%%%%%%%%%%%%%%%%%%%%%%%%%%%%%%%%%%%%%%%%%%%%%%%
\subsection{Binomial Versus Hypergeometric}
\label{subsec:binomial-versus-hypergeometric}
%%%%%%%%%%%%%%%%%%%%%%%%%%%%%%%%%%%%%%%%%%%%%%%%%%%%%%%%%%%%

The major distinction is replacement and independence.

\[
\begin{array}{c|c|c}
&
\text{Binomial}
&
\text{Hypergeometric}
\\
\hline
\text{Trials}
&
\text{fixed number}
&
\text{fixed sample size}
\\
\text{Replacement}
&
\text{effectively yes}
&
\text{no}
\\
\text{Independence}
&
\text{yes}
&
\text{no}
\\
\text{Success probability}
&
\text{constant}
&
\text{changes after draws}
\end{array}
\]

If sampling is without replacement from a finite population,
hypergeometric is usually the exact model.

%%%%%%%%%%%%%%%%%%%%%%%%%%%%%%%%%%%%%%%%%%%%%%%%%%%%%%%%%%%%
\subsection{Large Population Approximation}
\label{subsec:hypergeometric-binomial-approximation}
%%%%%%%%%%%%%%%%%%%%%%%%%%%%%%%%%%%%%%%%%%%%%%%%%%%%%%%%%%%%

If the population is very large relative to the sample, removing one
object changes the success proportion only slightly.

Then a hypergeometric variable may be approximated by

\[
\boxed{
\operatorname{Bin}\left(n,\frac KN\right).
}
\]

The approximation improves when the sampling fraction

\[
\frac nN
\]

is small.

%%%%%%%%%%%%%%%%%%%%%%%%%%%%%%%%%%%%%%%%%%%%%%%%%%%%%%%%%%%%
\subsection{Poisson Distribution}
\label{subsec:poisson-distribution}
%%%%%%%%%%%%%%%%%%%%%%%%%%%%%%%%%%%%%%%%%%%%%%%%%%%%%%%%%%%%

The Poisson distribution models counts of events occurring in a fixed
region or interval.

Examples include:

\begin{itemize}
    \item calls arriving in one minute;
    \item defects occurring along a length of material;
    \item particles detected during an interval;
    \item rare accidents during a period;
    \item customers arriving at a service station.
\end{itemize}

%%%%%%%%%%%%%%%%%%%%%%%%%%%%%%%%%%%%%%%%%%%%%%%%%%%%%%%%%%%%
\subsection{Poisson PMF}
\label{subsec:poisson-pmf}
%%%%%%%%%%%%%%%%%%%%%%%%%%%%%%%%%%%%%%%%%%%%%%%%%%%%%%%%%%%%

\begin{definitionbox}{Poisson Distribution}
A random variable \(X\) has a Poisson distribution with parameter
\(\lambda>0\), written

\[
\boxed{
X\sim\operatorname{Poisson}(\lambda),
}
\]

if

\[
\boxed{
P(X=k)
=
e^{-\lambda}
\frac{\lambda^k}{k!},
}
\]

for

\[
k=0,1,2,\ldots.
\]
\end{definitionbox}

The parameter

\[
\lambda
\]

is called ``lambda.''

It represents the average event count over the specified interval in
the standard Poisson model.

%%%%%%%%%%%%%%%%%%%%%%%%%%%%%%%%%%%%%%%%%%%%%%%%%%%%%%%%%%%%
\subsection{Poisson Normalization}
\label{subsec:poisson-normalization}
%%%%%%%%%%%%%%%%%%%%%%%%%%%%%%%%%%%%%%%%%%%%%%%%%%%%%%%%%%%%

Using the exponential series,

\[
e^\lambda
=
\sum_{k=0}^{\infty}
\frac{\lambda^k}{k!},
\]

we obtain

\[
\sum_{k=0}^{\infty}
e^{-\lambda}
\frac{\lambda^k}{k!}
=
e^{-\lambda}e^\lambda.
\]

Therefore,

\[
\boxed{
\sum_{k=0}^{\infty}P(X=k)=1.
}
\]

%%%%%%%%%%%%%%%%%%%%%%%%%%%%%%%%%%%%%%%%%%%%%%%%%%%%%%%%%%%%
\subsection{Worked Example: Poisson Arrivals}
\label{subsec:worked-poisson-arrivals}
%%%%%%%%%%%%%%%%%%%%%%%%%%%%%%%%%%%%%%%%%%%%%%%%%%%%%%%%%%%%

\begin{workedexamplebox}{Five Arrivals per Hour}

Suppose customers arrive according to a Poisson model with average

\[
5
\]

customers per hour.

Let \(X\) be the number arriving during one hour.

Then

\[
X\sim\operatorname{Poisson}(5).
\]

The probability of exactly \(3\) arrivals is

\[
P(X=3)
=
e^{-5}
\frac{5^3}{3!}.
\]

\begin{answerbox}
\[
\boxed{
P(X=3)
=
e^{-5}
\frac{125}{6}.
}
\]
\end{answerbox}

\end{workedexamplebox}

%%%%%%%%%%%%%%%%%%%%%%%%%%%%%%%%%%%%%%%%%%%%%%%%%%%%%%%%%%%%
\subsection{Scaling the Poisson Rate}
\label{subsec:poisson-rate-scaling}
%%%%%%%%%%%%%%%%%%%%%%%%%%%%%%%%%%%%%%%%%%%%%%%%%%%%%%%%%%%%

Suppose the average number of events is

\[
\lambda
\]

per unit time.

Over an interval of length \(t\), the expected count becomes

\[
\lambda t.
\]

Thus the count may be modeled by

\[
\boxed{
X_t
\sim
\operatorname{Poisson}(\lambda t).
}
\]

Care must be taken to keep the units consistent.

%%%%%%%%%%%%%%%%%%%%%%%%%%%%%%%%%%%%%%%%%%%%%%%%%%%%%%%%%%%%
\subsection{Worked Example: Changing the Time Interval}
\label{subsec:worked-poisson-scaling}
%%%%%%%%%%%%%%%%%%%%%%%%%%%%%%%%%%%%%%%%%%%%%%%%%%%%%%%%%%%%

\begin{workedexamplebox}{Ten Calls per Hour}

Suppose calls arrive at an average rate of

\[
10
\]

per hour.

Over \(30\) minutes,

\[
t=\frac12\text{ hour}.
\]

Therefore,

\[
\lambda t
=
10\left(\frac12\right)
=
5.
\]

The number \(X\) of calls in \(30\) minutes is modeled by

\[
\begin{answerbox}
\[
\boxed{
X\sim\operatorname{Poisson}(5).
}
\]
\end{answerbox}

\end{workedexamplebox}

%%%%%%%%%%%%%%%%%%%%%%%%%%%%%%%%%%%%%%%%%%%%%%%%%%%%%%%%%%%%
\subsection{Poisson Recurrence}
\label{subsec:poisson-recurrence}
%%%%%%%%%%%%%%%%%%%%%%%%%%%%%%%%%%%%%%%%%%%%%%%%%%%%%%%%%%%%

If

\[
X\sim\operatorname{Poisson}(\lambda),
\]

then

\[
P(X=k)
=
e^{-\lambda}
\frac{\lambda^k}{k!}.
\]

Therefore,

\[
\frac{P(X=k+1)}{P(X=k)}
=
\frac{\lambda}{k+1}.
\]

Hence,

\[
\boxed{
P(X=k+1)
=
P(X=k)
\frac{\lambda}{k+1}.
}
\]

%%%%%%%%%%%%%%%%%%%%%%%%%%%%%%%%%%%%%%%%%%%%%%%%%%%%%%%%%%%%
\subsection{Poisson Approximation to the Binomial}
\label{subsec:poisson-binomial-approximation}
%%%%%%%%%%%%%%%%%%%%%%%%%%%%%%%%%%%%%%%%%%%%%%%%%%%%%%%%%%%%

Suppose

\[
X\sim\operatorname{Bin}(n,p).
\]

When

\[
n
\]

is large,

\[
p
\]

is small, and

\[
np
\]

has moderate size, the binomial distribution can often be approximated
by

\[
\boxed{
Y\sim\operatorname{Poisson}(np).
}
\]

Thus,

\[
\boxed{
\lambda=np.
}
\]

%%%%%%%%%%%%%%%%%%%%%%%%%%%%%%%%%%%%%%%%%%%%%%%%%%%%%%%%%%%%
\subsection{Deriving the Poisson Limit}
\label{subsec:derive-poisson-limit}
%%%%%%%%%%%%%%%%%%%%%%%%%%%%%%%%%%%%%%%%%%%%%%%%%%%%%%%%%%%%

Let

\[
p=\frac{\lambda}{n}.
\]

Then

\[
P(X=k)
=
\binom nk
\left(\frac{\lambda}{n}\right)^k
\left(1-\frac{\lambda}{n}\right)^{n-k}.
\]

For fixed \(k\), as

\[
n\to\infty,
\]

we have

\[
\binom nk
\left(\frac1n\right)^k
\to
\frac1{k!},
\]

and

\[
\left(1-\frac{\lambda}{n}\right)^n
\to
e^{-\lambda}.
\]

Therefore,

\[
\boxed{
P(X=k)
\to
e^{-\lambda}
\frac{\lambda^k}{k!}.
}
\]

%%%%%%%%%%%%%%%%%%%%%%%%%%%%%%%%%%%%%%%%%%%%%%%%%%%%%%%%%%%%
\subsection{Worked Example: Rare Defects}
\label{subsec:worked-poisson-approx-binomial}
%%%%%%%%%%%%%%%%%%%%%%%%%%%%%%%%%%%%%%%%%%%%%%%%%%%%%%%%%%%%

\begin{workedexamplebox}{Rare Defects in a Large Batch}

Suppose

\[
X\sim\operatorname{Bin}(1000,0.002).
\]

Then

\[
np
=
1000(0.002)
=
2.
\]

A Poisson approximation uses

\[
Y\sim\operatorname{Poisson}(2).
\]

Thus,

\[
P(X=3)
\approx
P(Y=3)
=
e^{-2}\frac{2^3}{3!}.
\]

\begin{answerbox}
\[
\boxed{
P(X=3)
\approx
e^{-2}\frac{8}{6}.
}
\]
\end{answerbox}

\end{workedexamplebox}

%%%%%%%%%%%%%%%%%%%%%%%%%%%%%%%%%%%%%%%%%%%%%%%%%%%%%%%%%%%%
\subsection{Categorical Distribution}
\label{subsec:categorical-distribution}
%%%%%%%%%%%%%%%%%%%%%%%%%%%%%%%%%%%%%%%%%%%%%%%%%%%%%%%%%%%%

A categorical experiment has one outcome from several possible
categories.

Suppose the categories are

\[
1,\ldots,m
\]

with probabilities

\[
p_1,\ldots,p_m,
\]

where

\[
p_i\geq0
\]

and

\[
\sum_{i=1}^{m}p_i=1.
\]

A single categorical draw generalizes a Bernoulli trial.

%%%%%%%%%%%%%%%%%%%%%%%%%%%%%%%%%%%%%%%%%%%%%%%%%%%%%%%%%%%%
\subsection{Multinomial Distribution}
\label{subsec:multinomial-distribution}
%%%%%%%%%%%%%%%%%%%%%%%%%%%%%%%%%%%%%%%%%%%%%%%%%%%%%%%%%%%%

The multinomial distribution generalizes the binomial distribution to
more than two categories.

Suppose \(n\) independent trials each produce one of \(m\) categories
with probabilities

\[
p_1,\ldots,p_m.
\]

Let

\[
X_i
=
\text{number of outcomes in category }i.
\]

Then

\[
X_1+\cdots+X_m=n.
\]

%%%%%%%%%%%%%%%%%%%%%%%%%%%%%%%%%%%%%%%%%%%%%%%%%%%%%%%%%%%%
\subsection{Multinomial PMF}
\label{subsec:multinomial-pmf}
%%%%%%%%%%%%%%%%%%%%%%%%%%%%%%%%%%%%%%%%%%%%%%%%%%%%%%%%%%%%

For nonnegative integers

\[
x_1,\ldots,x_m
\]

satisfying

\[
x_1+\cdots+x_m=n,
\]

the multinomial probability is

\[
\boxed{
P(X_1=x_1,\ldots,X_m=x_m)
=
\frac{n!}{x_1!\cdots x_m!}
p_1^{x_1}\cdots p_m^{x_m}.
}
\]

The coefficient

\[
\frac{n!}{x_1!\cdots x_m!}
\]

counts arrangements of the category labels.

%%%%%%%%%%%%%%%%%%%%%%%%%%%%%%%%%%%%%%%%%%%%%%%%%%%%%%%%%%%%
\subsection{Worked Example: Multinomial Counts}
\label{subsec:worked-multinomial}
%%%%%%%%%%%%%%%%%%%%%%%%%%%%%%%%%%%%%%%%%%%%%%%%%%%%%%%%%%%%

\begin{workedexamplebox}{Three Categories}

Suppose each trial produces category \(A\), \(B\), or \(C\) with

\[
P(A)=0.5,
\qquad
P(B)=0.3,
\qquad
P(C)=0.2.
\]

Perform \(10\) independent trials.

Find the probability of exactly

\[
5
\]

results of type \(A\),

\[
3
\]

of type \(B\), and

\[
2
\]

of type \(C\).

\begin{answerbox}
\[
\boxed{
P
=
\frac{10!}{5!3!2!}
(0.5)^5
(0.3)^3
(0.2)^2.
}
\]
\end{answerbox}

\end{workedexamplebox}

%%%%%%%%%%%%%%%%%%%%%%%%%%%%%%%%%%%%%%%%%%%%%%%%%%%%%%%%%%%%
\subsection{Binomial as a Multinomial Special Case}
\label{subsec:binomial-multinomial-special}
%%%%%%%%%%%%%%%%%%%%%%%%%%%%%%%%%%%%%%%%%%%%%%%%%%%%%%%%%%%%

When

\[
m=2,
\]

the multinomial PMF becomes

\[
\frac{n!}{k!(n-k)!}
p^k(1-p)^{n-k}.
\]

Thus,

\[
\boxed{
\text{binomial}
=
\text{two-category multinomial}.
}
\]

%%%%%%%%%%%%%%%%%%%%%%%%%%%%%%%%%%%%%%%%%%%%%%%%%%%%%%%%%%%%
\subsection{Poisson Binomial Distribution}
\label{subsec:poisson-binomial}
%%%%%%%%%%%%%%%%%%%%%%%%%%%%%%%%%%%%%%%%%%%%%%%%%%%%%%%%%%%%

Suppose

\[
X_1,\ldots,X_n
\]

are independent Bernoulli random variables with possibly different
success probabilities

\[
p_1,\ldots,p_n.
\]

Then

\[
X
=
X_1+\cdots+X_n
\]

has a Poisson binomial distribution.

If

\[
p_1=\cdots=p_n=p,
\]

this reduces to the ordinary binomial distribution.

%%%%%%%%%%%%%%%%%%%%%%%%%%%%%%%%%%%%%%%%%%%%%%%%%%%%%%%%%%%%
\subsection{Worked Example: Unequal Bernoulli Probabilities}
\label{subsec:worked-poisson-binomial}
%%%%%%%%%%%%%%%%%%%%%%%%%%%%%%%%%%%%%%%%%%%%%%%%%%%%%%%%%%%%

\begin{workedexamplebox}{Two Unequal Trials}

Let

\[
X_1\sim\operatorname{Bernoulli}(0.2)
\]

and

\[
X_2\sim\operatorname{Bernoulli}(0.7),
\]

independently.

Let

\[
X=X_1+X_2.
\]

Then

\[
P(X=0)
=
(0.8)(0.3)
=
0.24,
\]

\[
P(X=2)
=
(0.2)(0.7)
=
0.14.
\]

Hence,

\[
P(X=1)
=
1-0.24-0.14
=
0.62.
\]

\begin{answerbox}
\[
\boxed{
P(X=0)=0.24,\quad
P(X=1)=0.62,\quad
P(X=2)=0.14.
}
\]
\end{answerbox}

\end{workedexamplebox}

%%%%%%%%%%%%%%%%%%%%%%%%%%%%%%%%%%%%%%%%%%%%%%%%%%%%%%%%%%%%
\subsection{Discrete Waiting-Time Models}
\label{subsec:discrete-waiting-time-models}
%%%%%%%%%%%%%%%%%%%%%%%%%%%%%%%%%%%%%%%%%%%%%%%%%%%%%%%%%%%%

Several discrete distributions model waiting.

\[
\boxed{
\operatorname{Geom}(p)
}
\]

models waiting until the first success.

\[
\boxed{
\operatorname{NegBin}(r,p)
}
\]

models waiting until the \(r\)-th success.

The support depends on whether the variable counts:

\[
\text{trials},
\]

or

\[
\text{failures}.
\]

The convention must always be stated.

%%%%%%%%%%%%%%%%%%%%%%%%%%%%%%%%%%%%%%%%%%%%%%%%%%%%%%%%%%%%
\subsection{Discrete Count Models}
\label{subsec:discrete-count-models}
%%%%%%%%%%%%%%%%%%%%%%%%%%%%%%%%%%%%%%%%%%%%%%%%%%%%%%%%%%%%

Common count distributions include:

\[
\boxed{
\operatorname{Bin}(n,p)
}
\]

for successes in a fixed number of independent trials,

\[
\boxed{
\operatorname{Hypergeom}(N,K,n)
}
\]

for successes in sampling without replacement,

and

\[
\boxed{
\operatorname{Poisson}(\lambda)
}
\]

for events occurring over a fixed interval.

Choosing the correct count model depends primarily on the data-generating
mechanism.

%%%%%%%%%%%%%%%%%%%%%%%%%%%%%%%%%%%%%%%%%%%%%%%%%%%%%%%%%%%%
\subsection{Distribution Selection Table}
\label{subsec:discrete-distribution-selection-table}
%%%%%%%%%%%%%%%%%%%%%%%%%%%%%%%%%%%%%%%%%%%%%%%%%%%%%%%%%%%%

\[
\begin{array}{c|c}
\text{Problem Structure}
&
\text{Distribution}
\\
\hline
\text{One success/failure trial}
&
\operatorname{Bernoulli}(p)
\\[1mm]
\text{Successes in \(n\) independent trials}
&
\operatorname{Bin}(n,p)
\\[1mm]
\text{Trials until first success}
&
\operatorname{Geom}(p)
\\[1mm]
\text{Trials until \(r\)-th success}
&
\operatorname{NegBin}(r,p)
\\[1mm]
\text{Successes without replacement}
&
\operatorname{Hypergeom}(N,K,n)
\\[1mm]
\text{Event count in fixed interval}
&
\operatorname{Poisson}(\lambda)
\\[1mm]
\text{Several categories over \(n\) trials}
&
\operatorname{Multinomial}
\end{array}
\]

%%%%%%%%%%%%%%%%%%%%%%%%%%%%%%%%%%%%%%%%%%%%%%%%%%%%%%%%%%%%
\subsection{CDF of a Discrete Distribution}
\label{subsec:discrete-distribution-cdf}
%%%%%%%%%%%%%%%%%%%%%%%%%%%%%%%%%%%%%%%%%%%%%%%%%%%%%%%%%%%%

For a discrete random variable,

\[
F_X(x)
=
P(X\leq x)
=
\sum_{t\leq x}p_X(t).
\]

For integer-valued \(X\),

\[
\boxed{
F_X(k)
=
\sum_{j\leq k}P(X=j).
}
\]

Thus cumulative probabilities are obtained by summing the PMF.

%%%%%%%%%%%%%%%%%%%%%%%%%%%%%%%%%%%%%%%%%%%%%%%%%%%%%%%%%%%%
\subsection{Tail Probabilities}
\label{subsec:discrete-tail-probabilities}
%%%%%%%%%%%%%%%%%%%%%%%%%%%%%%%%%%%%%%%%%%%%%%%%%%%%%%%%%%%%

The upper tail is

\[
P(X\geq k).
\]

For integer-valued variables,

\[
\boxed{
P(X\geq k)
=
1-P(X\leq k-1).
}
\]

Similarly,

\[
\boxed{
P(X>k)
=
1-P(X\leq k).
}
\]

This distinction matters for discrete random variables because point
probabilities may be positive.

%%%%%%%%%%%%%%%%%%%%%%%%%%%%%%%%%%%%%%%%%%%%%%%%%%%%%%%%%%%%
\subsection{Survival Function}
\label{subsec:discrete-survival-function}
%%%%%%%%%%%%%%%%%%%%%%%%%%%%%%%%%%%%%%%%%%%%%%%%%%%%%%%%%%%%

The survival function is

\[
\boxed{
S_X(x)
=
P(X>x).
}
\]

Since

\[
F_X(x)=P(X\leq x),
\]

we have

\[
\boxed{
S_X(x)
=
1-F_X(x).
}
\]

For waiting-time variables, survival probabilities are often easier to
work with than the PMF.

%%%%%%%%%%%%%%%%%%%%%%%%%%%%%%%%%%%%%%%%%%%%%%%%%%%%%%%%%%%%
\subsection{Probability Generating Functions Preview}
\label{subsec:discrete-pgf-preview}
%%%%%%%%%%%%%%%%%%%%%%%%%%%%%%%%%%%%%%%%%%%%%%%%%%%%%%%%%%%%

For a nonnegative integer-valued random variable \(X\), its probability
generating function is

\[
\boxed{
G_X(s)
=
\mathbb{E}[s^X].
}
\]

Expectation is developed systematically in the next several sections,
but formally this becomes

\[
\boxed{
G_X(s)
=
\sum_{k=0}^{\infty}
P(X=k)s^k.
}
\]

The coefficients of the power series are the probabilities.

%%%%%%%%%%%%%%%%%%%%%%%%%%%%%%%%%%%%%%%%%%%%%%%%%%%%%%%%%%%%
\subsection{PGF of a Bernoulli Variable}
\label{subsec:bernoulli-pgf}
%%%%%%%%%%%%%%%%%%%%%%%%%%%%%%%%%%%%%%%%%%%%%%%%%%%%%%%%%%%%

If

\[
X\sim\operatorname{Bernoulli}(p),
\]

then

\[
G_X(s)
=
(1-p)s^0+ps^1.
\]

Therefore,

\[
\boxed{
G_X(s)
=
1-p+ps.
}
\]

%%%%%%%%%%%%%%%%%%%%%%%%%%%%%%%%%%%%%%%%%%%%%%%%%%%%%%%%%%%%
\subsection{PGF of a Binomial Variable}
\label{subsec:binomial-pgf}
%%%%%%%%%%%%%%%%%%%%%%%%%%%%%%%%%%%%%%%%%%%%%%%%%%%%%%%%%%%%

For

\[
X\sim\operatorname{Bin}(n,p),
\]

\[
G_X(s)
=
\sum_{k=0}^{n}
\binom nk
p^k(1-p)^{n-k}s^k.
\]

By the binomial theorem,

\[
\boxed{
G_X(s)
=
(1-p+ps)^n.
}
\]

This mirrors the generating-function methods developed in Chapter 6.

%%%%%%%%%%%%%%%%%%%%%%%%%%%%%%%%%%%%%%%%%%%%%%%%%%%%%%%%%%%%
\subsection{PGF of a Poisson Variable}
\label{subsec:poisson-pgf}
%%%%%%%%%%%%%%%%%%%%%%%%%%%%%%%%%%%%%%%%%%%%%%%%%%%%%%%%%%%%

If

\[
X\sim\operatorname{Poisson}(\lambda),
\]

then

\[
G_X(s)
=
\sum_{k=0}^{\infty}
e^{-\lambda}
\frac{\lambda^k}{k!}
s^k.
\]

Thus,

\[
G_X(s)
=
e^{-\lambda}
\sum_{k=0}^{\infty}
\frac{(\lambda s)^k}{k!}.
\]

Therefore,

\[
\boxed{
G_X(s)
=
e^{\lambda(s-1)}.
}
\]

%%%%%%%%%%%%%%%%%%%%%%%%%%%%%%%%%%%%%%%%%%%%%%%%%%%%%%%%%%%%
\subsection{Independent Sums and PGFs Preview}
\label{subsec:independent-sums-pgf-preview}
%%%%%%%%%%%%%%%%%%%%%%%%%%%%%%%%%%%%%%%%%%%%%%%%%%%%%%%%%%%%

If \(X\) and \(Y\) are independent nonnegative integer-valued random
variables, then

\[
\boxed{
G_{X+Y}(s)
=
G_X(s)G_Y(s).
}
\]

This mirrors convolution of discrete distributions.

It also explains several closure properties of standard distributions.

%%%%%%%%%%%%%%%%%%%%%%%%%%%%%%%%%%%%%%%%%%%%%%%%%%%%%%%%%%%%
\subsection{Sum of Independent Bernoulli Variables}
\label{subsec:sum-independent-bernoulli}
%%%%%%%%%%%%%%%%%%%%%%%%%%%%%%%%%%%%%%%%%%%%%%%%%%%%%%%%%%%%

If

\[
X_i\sim\operatorname{Bernoulli}(p)
\]

independently for

\[
i=1,\ldots,n,
\]

then

\[
X
=
X_1+\cdots+X_n
\]

counts the total number of successes.

Therefore,

\[
\boxed{
X\sim\operatorname{Bin}(n,p).
}
\]

This is the structural origin of the binomial distribution.

%%%%%%%%%%%%%%%%%%%%%%%%%%%%%%%%%%%%%%%%%%%%%%%%%%%%%%%%%%%%
\subsection{Sum of Independent Poisson Variables}
\label{subsec:sum-independent-poisson}
%%%%%%%%%%%%%%%%%%%%%%%%%%%%%%%%%%%%%%%%%%%%%%%%%%%%%%%%%%%%

Suppose

\[
X\sim\operatorname{Poisson}(\lambda_1)
\]

and

\[
Y\sim\operatorname{Poisson}(\lambda_2)
\]

independently.

Then

\[
\boxed{
X+Y
\sim
\operatorname{Poisson}(\lambda_1+\lambda_2).
}
\]

This closure property is particularly important for combining counts
from disjoint time intervals or regions.

%%%%%%%%%%%%%%%%%%%%%%%%%%%%%%%%%%%%%%%%%%%%%%%%%%%%%%%%%%%%
\subsection{Worked Example: Combining Poisson Counts}
\label{subsec:worked-poisson-sum}
%%%%%%%%%%%%%%%%%%%%%%%%%%%%%%%%%%%%%%%%%%%%%%%%%%%%%%%%%%%%

\begin{workedexamplebox}{Two Independent Arrival Streams}

Suppose one arrival stream produces

\[
X\sim\operatorname{Poisson}(3)
\]

events per hour and another independent stream produces

\[
Y\sim\operatorname{Poisson}(5).
\]

Then the total count

\[
T=X+Y
\]

satisfies

\begin{answerbox}
\[
\boxed{
T\sim\operatorname{Poisson}(8).
}
\]
\end{answerbox}

\end{workedexamplebox}

%%%%%%%%%%%%%%%%%%%%%%%%%%%%%%%%%%%%%%%%%%%%%%%%%%%%%%%%%%%%
\subsection{Poisson Splitting Preview}
\label{subsec:poisson-splitting-preview}
%%%%%%%%%%%%%%%%%%%%%%%%%%%%%%%%%%%%%%%%%%%%%%%%%%%%%%%%%%%%

Suppose events occur according to a Poisson count with parameter
\(\lambda\).

Independently classify each event as type \(A\) with probability \(p\)
and type \(B\) with probability \(1-p\).

Then the type counts behave as independent Poisson variables with
parameters

\[
\boxed{
\lambda p
}
\]

and

\[
\boxed{
\lambda(1-p).
}
\]

This is called Poisson thinning or splitting.

%%%%%%%%%%%%%%%%%%%%%%%%%%%%%%%%%%%%%%%%%%%%%%%%%%%%%%%%%%%%
\subsection{Occupancy Problems}
\label{subsec:occupancy-distributions}
%%%%%%%%%%%%%%%%%%%%%%%%%%%%%%%%%%%%%%%%%%%%%%%%%%%%%%%%%%%%

An occupancy problem places objects into boxes.

For example:

\[
n
\]

balls may be independently assigned to

\[
m
\]

boxes.

Random variables of interest include:

\begin{itemize}
    \item number of empty boxes;
    \item number of boxes containing exactly one ball;
    \item maximum occupancy;
    \item number of collisions.
\end{itemize}

These variables need not follow one of the standard distributions
exactly, but binomial, multinomial, and indicator methods are often
useful.

%%%%%%%%%%%%%%%%%%%%%%%%%%%%%%%%%%%%%%%%%%%%%%%%%%%%%%%%%%%%
\subsection{Worked Example: Occupancy of One Box}
\label{subsec:worked-occupancy-one-box}
%%%%%%%%%%%%%%%%%%%%%%%%%%%%%%%%%%%%%%%%%%%%%%%%%%%%%%%%%%%%

\begin{workedexamplebox}{Balls into Boxes}

Place \(n\) balls independently and uniformly into \(m\) boxes.

Fix one particular box.

Each ball enters that box with probability

\[
\frac1m.
\]

Therefore, if \(X\) is the number of balls in that box,

\begin{answerbox}
\[
\boxed{
X
\sim
\operatorname{Bin}
\left(
n,\frac1m
\right).
}
\]
\end{answerbox}

\end{workedexamplebox}

%%%%%%%%%%%%%%%%%%%%%%%%%%%%%%%%%%%%%%%%%%%%%%%%%%%%%%%%%%%%
\subsection{Coupon Collector Preview}
\label{subsec:coupon-collector-preview}
%%%%%%%%%%%%%%%%%%%%%%%%%%%%%%%%%%%%%%%%%%%%%%%%%%%%%%%%%%%%

Suppose each trial independently produces one of \(m\) coupon types,
uniformly.

Let

\[
T
\]

be the number of trials needed to observe every type.

The distribution of \(T\) is not geometric because the success
probability changes as new coupon types are collected.

However, the waiting time can be decomposed into stages that resemble
geometric random variables.

This problem becomes useful later when studying expectation.

%%%%%%%%%%%%%%%%%%%%%%%%%%%%%%%%%%%%%%%%%%%%%%%%%%%%%%%%%%%%
\subsection{Discrete Reliability Models}
\label{subsec:discrete-reliability-models}
%%%%%%%%%%%%%%%%%%%%%%%%%%%%%%%%%%%%%%%%%%%%%%%%%%%%%%%%%%%%

Suppose \(n\) independent components each function with probability

\[
p.
\]

Then the number of functioning components is

\[
\boxed{
X\sim\operatorname{Bin}(n,p).
}
\]

A system requiring at least \(r\) functioning components has reliability

\[
\boxed{
P(X\geq r)
=
\sum_{k=r}^{n}
\binom nk
p^k(1-p)^{n-k}.
}
\]

%%%%%%%%%%%%%%%%%%%%%%%%%%%%%%%%%%%%%%%%%%%%%%%%%%%%%%%%%%%%
\subsection{Worked Example: \(k\)-out-of-\(n\) System}
\label{subsec:worked-k-out-of-n}
%%%%%%%%%%%%%%%%%%%%%%%%%%%%%%%%%%%%%%%%%%%%%%%%%%%%%%%%%%%%

\begin{workedexamplebox}{System Reliability}

Suppose a system contains \(5\) independent components.

Each works with probability

\[
0.9.
\]

The system works if at least \(4\) components work.

Let

\[
X\sim\operatorname{Bin}(5,0.9).
\]

Then

\[
P(\text{system works})
=
P(X\geq4).
\]

Therefore,

\begin{answerbox}
\[
\boxed{
P(\text{system works})
=
\binom54
(0.9)^4(0.1)
+
(0.9)^5.
}
\]
\end{answerbox}

\end{workedexamplebox}

%%%%%%%%%%%%%%%%%%%%%%%%%%%%%%%%%%%%%%%%%%%%%%%%%%%%%%%%%%%%
\subsection{Discrete Distribution Relationships}
\label{subsec:discrete-distribution-relationships}
%%%%%%%%%%%%%%%%%%%%%%%%%%%%%%%%%%%%%%%%%%%%%%%%%%%%%%%%%%%%

Several standard distributions are connected.

\[
\boxed{
\operatorname{Bernoulli}(p)
=
\operatorname{Bin}(1,p).
}
\]

Also,

\[
\boxed{
\operatorname{Geom}(p)
=
\operatorname{NegBin}(1,p)
}
\]

under the convention used in this text.

For rare successes,

\[
\boxed{
\operatorname{Bin}(n,p)
\approx
\operatorname{Poisson}(np)
}
\]

when \(n\) is large and \(p\) is small.

For a large population with small sampling fraction,

\[
\boxed{
\operatorname{Hypergeom}(N,K,n)
\approx
\operatorname{Bin}
\left(
n,\frac KN
\right).
}
\]

%%%%%%%%%%%%%%%%%%%%%%%%%%%%%%%%%%%%%%%%%%%%%%%%%%%%%%%%%%%%
\subsection{Model Selection Workflow}
\label{subsec:discrete-model-selection-workflow}
%%%%%%%%%%%%%%%%%%%%%%%%%%%%%%%%%%%%%%%%%%%%%%%%%%%%%%%%%%%%

When choosing a discrete distribution, ask:

\begin{enumerate}
    \item Is there one trial or many?
    \item Is the variable a count or a waiting time?
    \item Is the number of trials fixed?
    \item Are trials independent?
    \item Is the success probability constant?
    \item Is sampling performed without replacement?
    \item Is the variable counting events in time or space?
    \item Are there two categories or several?
\end{enumerate}

These questions usually identify the appropriate model.

%%%%%%%%%%%%%%%%%%%%%%%%%%%%%%%%%%%%%%%%%%%%%%%%%%%%%%%%%%%%
\subsection{Common Errors}
\label{subsec:discrete-distributions-common-errors}
%%%%%%%%%%%%%%%%%%%%%%%%%%%%%%%%%%%%%%%%%%%%%%%%%%%%%%%%%%%%

\subsubsection{Using Binomial for Sampling Without Replacement}

If sampling is from a finite population without replacement, the trials
are generally dependent.

The exact model is typically hypergeometric.

%%%%%%%%%%%%%%%%%%%%%%%%%%%%%%%%%%%%%%%%%%%%%%%%%%%%%%%%%%%%
\subsubsection{Ignoring the Fixed-\(n\) Requirement}

A binomial variable counts successes in a fixed number of trials.

A geometric or negative binomial variable instead has a random number
of trials.

%%%%%%%%%%%%%%%%%%%%%%%%%%%%%%%%%%%%%%%%%%%%%%%%%%%%%%%%%%%%
\subsubsection{Confusing Geometric Conventions}

Some texts count trials until success.

Others count failures before success.

Always inspect the support.

%%%%%%%%%%%%%%%%%%%%%%%%%%%%%%%%%%%%%%%%%%%%%%%%%%%%%%%%%%%%
\subsubsection{Confusing Binomial and Negative Binomial}

Binomial:

\[
\boxed{
\text{fixed number of trials, random successes}.
}
\]

Negative binomial:

\[
\boxed{
\text{fixed number of successes, random number of trials}.
}
\]

%%%%%%%%%%%%%%%%%%%%%%%%%%%%%%%%%%%%%%%%%%%%%%%%%%%%%%%%%%%%
\subsubsection{Using Poisson Without Scaling the Interval}

If the rate is given per hour but the interval is measured in minutes,
convert the interval before determining the Poisson parameter.

%%%%%%%%%%%%%%%%%%%%%%%%%%%%%%%%%%%%%%%%%%%%%%%%%%%%%%%%%%%%
\subsubsection{Forgetting the Hypergeometric Denominator}

For unordered samples without replacement, the total number of samples
is

\[
\boxed{
\binom Nn.
}
\]

%%%%%%%%%%%%%%%%%%%%%%%%%%%%%%%%%%%%%%%%%%%%%%%%%%%%%%%%%%%%
\subsubsection{Ignoring the Hypergeometric Support}

Values of \(k\) must be feasible given both success and failure counts.

Not every value between \(0\) and \(n\) is always possible.

%%%%%%%%%%%%%%%%%%%%%%%%%%%%%%%%%%%%%%%%%%%%%%%%%%%%%%%%%%%%
\subsubsection{Using Poisson Approximation When \(p\) Is Not Small}

The Poisson approximation to a binomial distribution is mainly intended
for relatively rare events.

%%%%%%%%%%%%%%%%%%%%%%%%%%%%%%%%%%%%%%%%%%%%%%%%%%%%%%%%%%%%
\subsubsection{Forgetting Independent Trial Assumptions}

Binomial and geometric models assume repeated independent trials with
constant success probability.

%%%%%%%%%%%%%%%%%%%%%%%%%%%%%%%%%%%%%%%%%%%%%%%%%%%%%%%%%%%%
\subsubsection{Confusing PMF and CDF}

The PMF gives

\[
P(X=k).
\]

The CDF gives

\[
P(X\leq k).
\]

For discrete variables, the CDF is obtained by summing the PMF.

%%%%%%%%%%%%%%%%%%%%%%%%%%%%%%%%%%%%%%%%%%%%%%%%%%%%%%%%%%%%
\subsection{Applications}
\label{subsec:discrete-distributions-applications}
%%%%%%%%%%%%%%%%%%%%%%%%%%%%%%%%%%%%%%%%%%%%%%%%%%%%%%%%%%%%

\begin{applicationbox}

\textbf{Quality Control.}
Binomial and hypergeometric models describe defective items and sample
inspection.

\medskip

\textbf{Reliability Engineering.}
Bernoulli and binomial variables model functioning and failed
components.

\medskip

\textbf{Queueing Systems.}
Poisson counts are fundamental models for arrival processes.

\medskip

\textbf{Epidemiology.}
Binomial models describe case counts under fixed population assumptions,
while Poisson models often describe event incidence over time.

\medskip

\textbf{Finance.}
Discrete default indicators and event counts can be modeled using
Bernoulli, binomial, or Poisson distributions.

\medskip

\textbf{Computer Science.}
Hashing, randomized algorithms, packet arrivals, collisions, and
occupancy problems use discrete probability models.

\medskip

\textbf{Survey Sampling.}
Sampling without replacement naturally produces hypergeometric
distributions.

\medskip

\textbf{Biology.}
Genetic inheritance and repeated experimental outcomes frequently use
binomial and multinomial models.

\medskip

\textbf{Operations Research.}
Demand counts, arrivals, reliability, and inventory events often use
Poisson and related count models.

\end{applicationbox}

%%%%%%%%%%%%%%%%%%%%%%%%%%%%%%%%%%%%%%%%%%%%%%%%%%%%%%%%%%%%
\subsection{Exercises}
\label{subsec:discrete-distributions-exercises}
%%%%%%%%%%%%%%%%%%%%%%%%%%%%%%%%%%%%%%%%%%%%%%%%%%%%%%%%%%%%

\begin{exercise}[1]
Define a discrete probability mass function.
\end{exercise}

\begin{exercise}[1]
State the two conditions required for a valid PMF.
\end{exercise}

\begin{exercise}[1]
Define a Bernoulli random variable.
\end{exercise}

\begin{exercise}[1]
Show that

\[
\operatorname{Bernoulli}(p)
=
\operatorname{Bin}(1,p).
\]
\end{exercise}

\begin{exercise}[1]
State the support of

\[
X\sim\operatorname{Bin}(n,p).
\]
\end{exercise}

\begin{exercise}[1]
State the PMF of a binomial random variable.
\end{exercise}

\begin{exercise}[1]
State the PMF of the geometric distribution using the trial-counting
convention.
\end{exercise}

\begin{exercise}[1]
State the PMF of the Poisson distribution.
\end{exercise}

\begin{exercise}[1]
State the PMF of the hypergeometric distribution.
\end{exercise}

\begin{exercise}[1]
Explain the main difference between binomial and hypergeometric
sampling.
\end{exercise}

\begin{exercise}[2]
Let

\[
X\sim\operatorname{Bin}(8,0.4).
\]

Find

\[
P(X=3).
\]
\end{exercise}

\begin{exercise}[2]
For the same random variable, find

\[
P(X=0).
\]
\end{exercise}

\begin{exercise}[2]
For the same random variable, express

\[
P(X\geq1)
\]

using a complement.
\end{exercise}

\begin{exercise}[2]
A fair coin is tossed \(12\) times. Find the probability of exactly
\(7\) heads.
\end{exercise}

\begin{exercise}[2]
A basketball player makes a free throw with probability \(0.8\).
Assuming independence, find the probability of making exactly \(8\) of
\(10\) free throws.
\end{exercise}

\begin{exercise}[2]
Suppose

\[
X\sim\operatorname{Geom}(0.25).
\]

Find

\[
P(X=4).
\]
\end{exercise}

\begin{exercise}[2]
For the same variable, find

\[
P(X>5).
\]
\end{exercise}

\begin{exercise}[2]
Suppose

\[
X\sim\operatorname{NegBin}(4,0.3).
\]

Find the probability that the fourth success occurs on trial \(9\).
\end{exercise}

\begin{exercise}[2]
A shipment contains \(30\) items, \(6\) of which are defective. Select
\(5\) without replacement. Find the probability of exactly \(2\)
defective items.
\end{exercise}

\begin{exercise}[2]
For the previous problem, identify \(N\), \(K\), and \(n\).
\end{exercise}

\begin{exercise}[2]
Suppose

\[
X\sim\operatorname{Poisson}(4).
\]

Find

\[
P(X=2).
\]
\end{exercise}

\begin{exercise}[2]
For the same variable, find

\[
P(X=0).
\]
\end{exercise}

\begin{exercise}[2]
A call center receives an average of \(12\) calls per hour. Under a
Poisson model, determine the parameter for a \(15\)-minute interval.
\end{exercise}

\begin{exercise}[2]
Using the previous parameter, find the probability of no calls in
\(15\) minutes.
\end{exercise}

\begin{exercise}[2]
Ten independent trials produce categories \(A,B,C\) with probabilities

\[
0.2,\quad0.5,\quad0.3.
\]

Find the probability of counts

\[
2,\quad5,\quad3.
\]
\end{exercise}

\begin{exercise}[3]
Derive the binomial PMF from Bernoulli trials.
\end{exercise}

\begin{exercise}[3]
Use the binomial theorem to prove that the binomial PMF sums to \(1\).
\end{exercise}

\begin{exercise}[3]
Derive the geometric PMF.
\end{exercise}

\begin{exercise}[3]
Use an infinite geometric series to verify normalization of the
geometric PMF.
\end{exercise}

\begin{exercise}[3]
Prove the memoryless property of the geometric distribution.
\end{exercise}

\begin{exercise}[3]
Derive the negative binomial PMF.
\end{exercise}

\begin{exercise}[3]
Show explicitly that the geometric distribution is the \(r=1\) case of
the negative binomial distribution.
\end{exercise}

\begin{exercise}[3]
Derive the hypergeometric PMF by counting samples.
\end{exercise}

\begin{exercise}[3]
Determine the exact support of

\[
\operatorname{Hypergeom}(N,K,n).
\]
\end{exercise}

\begin{exercise}[3]
Derive the Poisson normalization identity using the exponential series.
\end{exercise}

\begin{exercise}[3]
Derive the recurrence

\[
P(X=k+1)
=
P(X=k)\frac{\lambda}{k+1}
\]

for a Poisson variable.
\end{exercise}

\begin{exercise}[3]
Starting from

\[
X_n\sim
\operatorname{Bin}
\left(
n,\frac{\lambda}{n}
\right),
\]

derive the Poisson limit for fixed \(k\).
\end{exercise}

\begin{exercise}[3]
Explain why a hypergeometric distribution approaches a binomial model
when the sampling fraction is small.
\end{exercise}

\begin{exercise}[3]
Show that the multinomial PMF sums to \(1\) using the multinomial
theorem.
\end{exercise}

\begin{exercise}[3]
Derive the PGF of a Bernoulli random variable.
\end{exercise}

\begin{exercise}[3]
Derive the PGF of a binomial random variable.
\end{exercise}

\begin{exercise}[3]
Derive the PGF of a Poisson random variable.
\end{exercise}

\begin{exercise}[3]
Use PGFs to explain why the sum of independent Poisson random variables
is Poisson.
\end{exercise}

\begin{exercise}[3]
A system contains \(8\) independent components, each functioning with
probability \(0.95\). Find the probability that at least \(7\) function.
\end{exercise}

\begin{exercise}[3]
In a random occupancy model, \(20\) balls are placed independently into
\(5\) boxes. Find the distribution of the number of balls in one fixed
box.
\end{exercise}

\begin{exercise}[4]
Study the Poisson binomial distribution and compare it with the ordinary
binomial distribution.
\end{exercise}

\begin{exercise}[4]
Investigate the beta-binomial distribution as a model with random
success probability.
\end{exercise}

\begin{exercise}[4]
Study the negative hypergeometric distribution.
\end{exercise}

\begin{exercise}[4]
Investigate compound Poisson distributions.
\end{exercise}

\begin{exercise}[4]
Study zero-inflated discrete distributions used when observed zeros are
more frequent than standard models predict.
\end{exercise}

\begin{exercise}[4]
Investigate overdispersion and explain why a negative binomial model may
be preferred over a Poisson model.
\end{exercise}

\begin{exercise}[4]
Study Poisson thinning and prove the independence of the split counts.
\end{exercise}

\begin{exercise}[4]
Investigate the Poisson process and derive the Poisson count
distribution from process assumptions.
\end{exercise}

\begin{exercise}[4]
Study occupancy distributions for the number of empty boxes.
\end{exercise}

\begin{exercise}[4]
Investigate the coupon collector problem and its decomposition into
geometric waiting times.
\end{exercise}

%%%%%%%%%%%%%%%%%%%%%%%%%%%%%%%%%%%%%%%%%%%%%%%%%%%%%%%%%%%%
\subsection{Conceptual Summary}
\label{subsec:discrete-distributions-summary}
%%%%%%%%%%%%%%%%%%%%%%%%%%%%%%%%%%%%%%%%%%%%%%%%%%%%%%%%%%%%

A discrete random variable has PMF

\[
\boxed{
p_X(x)=P(X=x)
}
\]

with

\[
\boxed{
p_X(x)\geq0,
\qquad
\sum_xp_X(x)=1.
}
\]

A Bernoulli variable satisfies

\[
\boxed{
X\sim\operatorname{Bernoulli}(p),
}
\]

with

\[
P(X=1)=p,
\qquad
P(X=0)=1-p.
\]

A binomial variable counts successes in \(n\) independent Bernoulli
trials:

\[
\boxed{
P(X=k)
=
\binom nk
p^k(1-p)^{n-k}.
}
\]

A geometric variable counts trials until the first success:

\[
\boxed{
P(X=k)
=
(1-p)^{k-1}p.
}
\]

A negative binomial variable counts trials until the \(r\)-th success:

\[
\boxed{
P(X=k)
=
\binom{k-1}{r-1}
p^r(1-p)^{k-r}.
}
\]

A hypergeometric variable counts successes in sampling without
replacement:

\[
\boxed{
P(X=k)
=
\frac{
\binom Kk
\binom{N-K}{n-k}
}{
\binom Nn
}.
}
\]

A Poisson variable satisfies

\[
\boxed{
P(X=k)
=
e^{-\lambda}
\frac{\lambda^k}{k!}.
}
\]

For rare binomial events,

\[
\boxed{
\operatorname{Bin}(n,p)
\approx
\operatorname{Poisson}(np).
}
\]

The multinomial distribution generalizes the binomial model to several
categories:

\[
\boxed{
P(X_1=x_1,\ldots,X_m=x_m)
=
\frac{n!}{x_1!\cdots x_m!}
\prod_{i=1}^{m}p_i^{x_i}.
}
\]

The main modeling distinctions are

\[
\boxed{
\text{count versus waiting time},
}
\]

\[
\boxed{
\text{replacement versus no replacement},
}
\]

and

\[
\boxed{
\text{fixed trials versus fixed successes}.
}
\]

Standard discrete distributions therefore organize recurring
probabilistic mechanisms into reusable mathematical models.

%%%%%%%%%%%%%%%%%%%%%%%%%%%%%%%%%%%%%%%%%%%%%%%%%%%%%%%%%%%%
\subsection{Section Connections}
\label{subsec:discrete-distributions-connections}
%%%%%%%%%%%%%%%%%%%%%%%%%%%%%%%%%%%%%%%%%%%%%%%%%%%%%%%%%%%%

\chapterconnection{
Discrete Distributions builds directly on the random-variable framework
of Section~\ref{sec:random-variables}. Bernoulli, binomial, geometric,
negative binomial, hypergeometric, Poisson, and multinomial models
provide the principal finite and countable probability families used
throughout statistics and stochastic modeling. The next section,
Continuous Distributions, replaces probability masses by probability
densities over intervals. Joint Distributions develops dependence among
multiple random variables, while Expectation and Moments derive
numerical summaries such as means and variances for all of these
families. Functions of Random Variables and Probability Transforms later
explain convolution, generating functions, and distribution
transformations. Stochastic Processes extends Poisson counts and
Bernoulli trials into time-indexed random systems.
}

%%%%%%%%%%%%%%%%%%%%%%%%%%%%%%%%%%%%%%%%%%%%%%%%%%%%%%%%%%%%
% SECTION SOURCE TRACKING
%%%%%%%%%%%%%%%%%%%%%%%%%%%%%%%%%%%%%%%%%%%%%%%%%%%%%%%%%%%%

%------------------------------------------------------------
% 7.3 Discrete Distributions
%------------------------------------------------------------
%
% Source basis:
%   - Standard elementary probability distributions.
%   - Standard discrete distribution theory.
%   - Standard combinatorial probability.
%
% Used for:
%   - discrete uniform distribution;
%   - Bernoulli distribution;
%   - Bernoulli trials;
%   - binomial distribution;
%   - binomial PMF derivation;
%   - binomial cumulative probabilities;
%   - binomial probability recurrence;
%   - binomial mode;
%   - geometric distribution;
%   - geometric tail probabilities;
%   - memoryless property;
%   - geometric convention differences;
%   - negative binomial distribution;
%   - geometric as a negative-binomial special case;
%   - hypergeometric distribution;
%   - hypergeometric support;
%   - sampling without replacement;
%   - binomial versus hypergeometric;
%   - large-population approximation;
%   - Poisson distribution;
%   - Poisson rate scaling;
%   - Poisson recurrence;
%   - Poisson approximation to binomial;
%   - Poisson limit;
%   - categorical distribution;
%   - multinomial distribution;
%   - Poisson binomial distribution;
%   - occupancy problems;
%   - reliability models;
%   - probability generating functions preview;
%   - Bernoulli, binomial, and Poisson PGFs;
%   - independent sums;
%   - Poisson splitting preview;
%   - distribution relationships;
%   - model selection;
%   - applications and exercises.
%
% Notes:
%   - Expectation and variance formulas are intentionally developed in
%     Section 7.6 rather than duplicated here.
%   - Probability-generating functions are developed more systematically
%     in Section 7.8.
%   - Poisson-process assumptions are developed later under Stochastic
%     Processes.
%   - Continuous analogues begin in Section 7.4.
%
%%%%%%%%%%%%%%%%%%%%%%%%%%%%%%%%%%%%%%%%%%%%%%%%%%%%%%%%%%%%